\reportsection{Binding}
\label{sec:binding}
Binding is a key property of any vector commitment scheme, and different schemes have
different notions of it~\cite{CY26}: subvector position binding, its single-position
specialisation, and the weak variant that requires an honestly generated commitment. We
look at subvector position binding of a Merkle commitment scheme and show that it can be
reduced to the collision resistance of the underlying hash function. For a given hash
function, collision resistance gives us a bound that we can evaluate against the desired
security parameter. Specifically, we show that this bound falls short of common target
$\lambda$-bit collision resistance for single field elements in small fields. The size of
the digest output by a concrete hasher is inextricably linked to the underlying field we
operate over. This realisation led us to flag a bug in Arkworks in which the hasher
implementation would output a single field element (digest size) regardless of the size of
the underlying field.

\begin{definition}[Subvector position binding]\label{def:vc-binding}
$\VC$ is \emph{subvector position binding} if no efficient adversary outputs a commitment
$\comm$ with two accepted openings
$(\idxset_b,\MTMessageSubVector_b,\opening_b)$ that disagree at some
$i\in\idxset_0\cap\idxset_1$, except with probability $\negl$.  The commitment
has no honest-generation requirement.
\end{definition}

\begin{proposition}[Binding from collision resistance]\label{thm:binding-cr}
Let $\MT[\RO]$ be a Merkle commitment over a perfect binary tree of depth $d$ whose leaf
encoding is injective and whose leaf and internal-node calls are domain separated.  If
the shared $\digestbits$-bit oracle receives at most $\querybudget$ distinct
queries during the complete experiment, then
\[
  \Adv{bind}(\ThetaVC)
  \le \Pr[\text{oracle-trace collision}]
  \le \frac{\querybudget(\querybudget-1)}{2^{\digestbits+1}},
\]
which meets $2^{-\secpar}$ when $\querybudget\le2^{h}$ and
$\secpar+1+2h\le\digestbits$.
\status{Lean-backed}
\end{proposition}

The binary tree is the case Chiesa and Yogev prove directly, and we take their proof as
given~\cite{CY26}: either the adversary's query trace already contains a collision, or the
two accepted openings of the minimal disputed index yield one across the two verification
traces, bounded by a union bound over the at most $(d+1)^2$ query pairs.  This is the
straightforward form the book states before improving on it, and the form \cref{eq:rule}
consumes.  The reduction is machine-checked in Lean~4~\cite{ZLeanVC26} for a general Merkle
shape, closed for singleton openings and argued on paper for a general index set
(\cref{app:binding-generic}).

Now that we have seen the reduction to collision resistance, consider how a hasher is
implemented concretely. Every hasher in \textsf{ark-mt} satisfies one \code{MerkleHasher}
trait (\cref{app:hasher-trait}), and the current Poseidon2 implementation sets its digest
to one field element, for every field $\mathbb{F}$ and every target $\lambda$. Nothing in
the trait or the implementation reads $\lambda$.

Now suppose we set a target security parameter $\lambda=128$. BabyBear's modulus is
$p=2013265921$, a $31$-bit prime, and the conservative bound used throughout this report is
$b=\lfloor\log_2 p\rfloor=30$ bits per element, so a one-element digest publishes a root drawn
from a $30$-bit space. Applying \Cref{lem:birthday} at $\digestbits=30$, a collision search
succeeds once the query count passes roughly $2^{15.5}\approx46{,}000$, well short of a
$\lambda=128$ target.

The book states this result plainly for a $31$-bit field. We ran
the search only to confirm that the one-element preset sits in that regime, and it produced two
trees over different messages sharing one root in $44$ seconds on a commodity laptop,
averaged over three seeded runs (\cref{app:demo}). The collision is cross-shape, between an
$8$-leaf binary tree and an $8$-leaf $8$-ary tree on the same Poseidon2 instance, and it
still breaks \Cref{def:vc-binding} as stated because \code{hash\_leaf} and
\code{hash\_nodes} carry no layer, index, or arity tag (\cref{lem:geometry}). Working
backwards from $\lambda$ and $\mathbb{F}$, \cref{eq:rule} asks for
$D=\lceil257/30\rceil=9$ elements at $\lambda=128$ and $h=64$, so fixing the hasher means
widening \code{Digest} from \code{F} to a nine-element array and hashing into all nine on
every call, which we price in \Cref{sec:cost}.